\section{Введение и структура курса}

\subsection*{О данном учебном пособии}

Данное учебное пособие разработано для систематического освоения навыков программирования беспилотных летательных аппаратов на примере квадрокоптера «Геоскан Пионер». Курс построен таким образом, чтобы обеспечить плавное погружение в предметную область: от изучения синтаксических основ языка программирования до реализации сложных автономных полетных заданий.

Материал структурирован согласно классическому дидактическому принципу «от простого к сложному». На начальных этапах основное внимание уделяется фундаментальным конструкциям языка Lua, необходимым для понимания логики работы скриптов. Впоследствии, по мере накопления теоретической базы, фокус смещается на прикладные аспекты: работу с периферийным оборудованием (светодиодная индикация), обработку событий и алгоритмизацию автономного поведения дрона.

Такой подход позволяет сформировать у обучающегося целостное представление о взаимодействии программного кода и аппаратной части робототехнического комплекса.

\subsection*{Цели и задачи обучения}

Основной целью курса является формирование устойчивых навыков разработки программного обеспечения для управления БПЛА. В процессе обучения решаются следующие задачи:

\begin{itemize}
    \item \textbf{Освоение инструментария:} Изучение синтаксиса и семантики языка Lua, а также специфики его применения во встраиваемых системах.
    \item \textbf{Аппаратное взаимодействие:} Получение навыков управления бортовыми системами квадрокоптера, включая систему световой индикации и таймеры.
    \item \textbf{Событийно-ориентированное программирование:} Понимание принципов неблокирующего кода и обработки асинхронных событий, что является критически важным для робототехники.
    \item \textbf{Алгоритмическое мышление:} Развитие способности декомпозировать сложные полетные задачи на последовательность элементарных команд и состояний.
\end{itemize}

\subsection*{Структура учебного курса (Дорожная карта)}

Учебный курс разделен на логически завершенные модули, каждый из которых является фундаментом для последующего:

\begin{enumerate}
    \item \textbf{Модуль 1: Подготовка и безопасность.}
    Вводный раздел, посвященный технике безопасности при эксплуатации БПЛА, настройке программного окружения и подготовке оборудования к работе.

    \item \textbf{Модуль 2: Основы языка Lua.}
    Базовый курс программирования. Здесь рассматриваются типы данных, переменные, условные операторы, циклы и таблицы. Этот модуль необходим тем, кто ранее не сталкивался с Lua.

    \item \textbf{Модуль 3: Периферия и таймеры.}
    Переход к прикладным задачам. Изучение методов управления RGB-светодиодами (модуль Ledbar) и работы с системным временем. Вводится понятие таймеров как основы для создания динамических сценариев без блокировки основного потока выполнения.

    \item \textbf{Модуль 4: Автономные миссии и события.}
    Продвинутый уровень. Разбор архитектуры автопилота, работа с событийной моделью (Events) и программирование полетных заданий. Обучающиеся научатся создавать скрипты для автоматического взлета, навигации по точкам и выполнения фигур пилотажа.

    \item \textbf{Приложение.}
    Справочные материалы, содержащие описание API функций, список системных событий и полезные фрагменты кода (сниппеты) для часто встречающихся задач.
\end{enumerate}

\subsection*{Планируемые результаты обучения}

По завершении курса обучающийся приобретет компетенции, необходимые для самостоятельной разработки полетных сценариев. Вы научитесь:
\begin{itemize}
    \item Создавать эффективный и читаемый код на языке Lua, соответствующий стандартам разработки.
    \item Использовать возможности бортовой периферии для индикации статусов и режимов работы дрона.
    \item Проектировать архитектуру программ, устойчивую к ошибкам и способную корректно обрабатывать штатные и нештатные ситуации.
    \item Реализовывать автономные миссии различной сложности, от простого зависания до пролета по сложному маршруту с выполнением полезной нагрузки.
\end{itemize}
